\chapter{Минимальные требования к новой модели спектрального рождения}
\label{ch:v6-spectral-transition-new-model-minimal-requirements-gate}

\section{Почему нужен контракт до выбора модели}

Итоговый статус версии VI показал: классификационный носитель и язык
перехода построены, но общий физический процесс не замкнут. Поэтому нельзя
сначала выбрать привлекательный портал, новую струну или дополнительное
поле, а затем подгонять критерии успеха. Сначала задаётся обязательный
контракт, одинаковый для всех кандидатов.

Контракт является необходимым, но не достаточным: прохождение пунктов не
гарантирует истинность модели, однако провал любого обязательного пункта
запрещает объявлять рождение материи выведенным.

\section{R0: типизированный физический носитель}

Новая модель обязана сохранять
\begin{equation}
 G_{\rm SM}=SU(3)_c\times SU(2)_L\times U(1)_Y,
 \qquad
 \mathcal H_{15}=\mathcal H_{12}\oplus\mathcal H_3,
\end{equation}
а также Real-структуру, градуировку и уже доказанные KO6/Toeplitz-классы.
Путь $D(s)$ или поле перехода должно быть определено на полном физическом
носителе и интертвинить правильные представления. Запрещено заменять
$\mathcal H_{15}$ пятнадцатью неразличимыми копиями.

\section{R1: один родительский функционал}

Должен существовать один gauge- и Real-инвариантный функционал
\begin{equation}
 S_{\rm parent}[\Phi,A,\Psi,\ldots],
\end{equation}
из которого следуют вакуум, переходная конфигурация и endpoint. Для всех
трёх используются одна мера, одна следовая нормировка и один набор
контрчленов. Портал, внесённый только в переходное уравнение, тест не
проходит.

\section{R2: эндогенный запуск}

Переход должен запускаться самим родителем одним из двух способов:
\begin{enumerate}
 \item физическая отрицательная мода гессиана после удаления gauge- и
 переносных нулевых мод;
 \item конечное евклидово bounce-седло с одной отрицательной модой.
\end{enumerate}
Во втором случае экспонента и квантовые поправки должны вычисляться из
того же действия \cite{new-model-coleman1977,new-model-callan-coleman1977}.
Внешний clock или quench допустим только как динамическая подсистема с
явным энергетическим балансом, а не как заданная функция времени.

\section{R3: вероятность и скорость}

Модель должна задавать нормированную меру траекторий или начальных
состояний и вычислимую скорость
\begin{equation}
 \Gamma=A\exp(-B),
\end{equation}
где $B$ является действием выведенного седла, а $A$ содержит определитель,
нулевые моды и физический временной масштаб. Если процесс термический,
температура и состояние среды также должны быть частью постановки.
Индексы, ранги и следовые веса не заменяют $A$ и $B$.

\section{R4: устойчивый локализованный endpoint}

Конечный объект должен быть решением полных нелинейных уравнений с конечной
энергией, ненулевой локализационной длиной и требуемыми зарядами. После
факторизации симметрий его физический гессиан должен быть неотрицателен,
а внутренний сектор --- иметь положительную щель либо доказанный механизм
метастабильности. Сфалерон не проходит этот тест, поскольку является
седлом \cite{new-model-klinkhamer-manton1984}.

\section{R5: слепое численное предсказание}

До сравнения с наблюдением модель должна выдать хотя бы один observable,
не равный новому свободному параметру:
\begin{equation}
 m_*,\qquad \Gamma_*,\qquad \tau_*,\qquad
 \frac{\Gamma_q}{\Gamma_\ell},\qquad \xi_*.
\end{equation}
Допустима безразмерная комбинация, если все входы уже фиксированы прежними
гейтами. Простое воспроизведение $15$, $1/7$ или $3:1$ не считается новым
предсказанием.

\section{R6: фальсифицируемость и сертификат провала}

Каждый пункт обязан иметь заранее объявленный стоп-критерий:
неверный intertwiner, отсутствие общего действия, отсутствие физической
отрицательной моды или bounce, расходящийся prefactor, отрицательная мода
endpoint или зависимость observable от нового свободного параметра.

\begin{theorem}[Контракт допуска новой модели]
Кандидат допускается к ветви спектрального рождения только при совместном
выполнении R0--R6. Успех одного классификационного, топологического или
солитонного подпункта не компенсирует провал динамики, скорости или
endpoint.
\end{theorem}

\section{Проверка текущей версии по контракту}

Версия VI проходит R0 лишь в части носителя и классов. Она не содержит
нового общего $S_{\rm parent}$, эндогенного trigger, вычислимой скорости,
устойчивого фермионного endpoint и слепого численного предсказания.
Следовательно, текущая архитектура не является допущенной моделью
физического рождения, хотя остаётся строгим классификационно-кинематическим
основанием для следующей версии.

\begin{center}
\begin{tabular}{p{0.11\textwidth}p{0.51\textwidth}p{0.22\textwidth}}
\toprule
Тест & Требование & Версия VI \\
\midrule
R0 & физический Real/gauge-носитель & частично пройден \\
R1 & один функционал & не пройден \\
R2 & эндогенный запуск & не пройден \\
R3 & мера и скорость & не пройден \\
R4 & устойчивый endpoint & не пройден \\
R5 & новое число & не пройден \\
R6 & стоп-сертификаты & пройден методически \\
\bottomrule
\end{tabular}
\end{center}

\begin{nogocriterion}
Автоматически отклоняются пять обходов: ручная вставка ранга, скрытый
портал, внешний quench без динамического источника, чтение топологии как
вероятности и объявление неустойчивого седла частицей.
\end{nogocriterion}

\section{Следующий гейт}

Следующий гейт
\path{version6_spectral_transition_new_model_candidate_menu_gate}
должен сравнить минимальные архитектуры только по R0--R6 и выбрать не
самую красивую, а минимальную не отвергнутую. Нулевая гипотеза: ни один
кандидат из прежних скрытых порталов контракт не проходит. Стоп-критерий:
если все кандидаты требуют независимых новых параметров без нового
observable, ветвь версии VI завершается без выбора модели.

\begin{protocol}
\begin{enumerate}
 \item обязательных тестов: \textbf{7};
 \item носитель и типизация: \textbf{R0};
 \item общий функционал: \textbf{R1};
 \item эндогенный запуск: \textbf{R2};
 \item вероятность и скорость: \textbf{R3};
 \item устойчивый endpoint: \textbf{R4};
 \item слепое численное предсказание: \textbf{R5};
 \item сертификаты провала: \textbf{R6};
 \item допуск версии VI: \textbf{нет};
 \item статус контракта: \textbf{необходимый, не достаточный}.
\end{enumerate}
\end{protocol}

Машинный сертификат сохранён в
\path{s2t_v6_spectral_transition_new_model_minimal_requirements_gate_results.json}.

\begin{thebibliography}{9}
\bibitem{new-model-coleman1977}
S.~Coleman, \emph{The Fate of the False Vacuum. I. Semiclassical Theory},
Physical Review D 15 (1977), 2929--2936.

\bibitem{new-model-callan-coleman1977}
C.~G.~Callan, S.~Coleman,
\emph{The Fate of the False Vacuum. II. First Quantum Corrections},
Physical Review D 16 (1977), 1762--1768.

\bibitem{new-model-klinkhamer-manton1984}
F.~R.~Klinkhamer, N.~S.~Manton,
\emph{A Saddle-Point Solution in the Weinberg--Salam Theory},
Physical Review D 30 (1984), 2212--2220.
\end{thebibliography}